% 论文实验逐条归纳（中文）— 每条实验独立成节，格式仿 MACPO 电力调度示例
\documentclass[11pt,a4paper]{ctexart}
\usepackage{geometry}
\usepackage{amsmath,amssymb}
\usepackage{booktabs}
\usepackage{longtable}
\usepackage{array}
\usepackage{graphicx}
\usepackage{hyperref}
\usepackage{xcolor}
\usepackage{enumitem}
\usepackage{algorithm}
\usepackage{algpseudocode}
\usepackage{etoolbox}
\geometry{margin=2cm}
\hypersetup{colorlinks=true,linkcolor=blue}
\setlist{nosep,leftmargin=*}

% 伪代码块后附中文变量说明（伪代码本身仅英文）
\newcommand{\varnote}[1]{\par\smallskip\noindent\textbf{变量说明（中文）：}#1\par\smallskip}

% 若图文件存在则嵌入，否则显示路径提示
\newcommand{\tryfig}[3][0.92\textwidth]{%
  \IfFileExists{#2}{%
    \begin{figure}[htbp]\centering
    \includegraphics[width=#1]{#2}
    \caption{#3}\end{figure}%
  }{%
    \begin{center}\fbox{\parbox{0.92\textwidth}{\centering
    \textbf{【实验结果图】} #3\\[0.3em]
    \small 图文件：\texttt{#2}（运行对应出图脚本后重新编译即可嵌入）%
    }}\end{center}%
  }%
}

\title{\textbf{RL-MACPO 论文实验归纳（中文阅读版）}\\[0.3em]
\large 含背景、术语表、公式出处 · 26 条实验逐条说明}
\author{对应英文稿 \texttt{conference\_new\_ready.tex} · 2026-07-10}
\date{}

\begin{document}
\maketitle

\begin{abstract}
本文档是对论文 \texttt{conference\_new\_ready.tex}（冲突门控通信 RL-MACPO）中\textbf{全部 26 项实验}的中文归纳，供读论文时对照查阅。\textbf{请先阅读第\ref{sec:background}--\ref{sec:varindex}节}（背景、术语表、变量速查、公式来源），再进入具体实验。伪代码块内仅英文变量名；中文含义见 §2.2 变量速查或各算法后的「变量说明」。每条实验结构：实验目的 $\to$ 公式 $\to$ 流程 $\to$ 指标 $\to$ 设置 $\to$ 结果。
\end{abstract}

\tableofcontents
\newpage

%==============================================================================
\section{背景：这篇论文在做什么}\label{sec:background}
%==============================================================================
\subsection{问题}
\textbf{网络分布式黑盒优化（NDO）}：$n$ 个智能体在图结构上各自优化局部目标，但相邻智能体共享一部分决策变量（如电网联络线功率、交通流共享路段）。经典算法 \textbf{MACPO}（Multi-Agent Cooperative Penalty Optimization，见 MACPO 原文） 的做法是：每轮外循环都\textbf{通信并协商}共享变量，用惩罚项把不一致性压下去。

\textbf{痛点}：很多外循环里冲突其实很小，每轮都协商浪费通信与计算，对最终解质量帮助不大。

\textbf{本文核心想法}：在每个外循环先问一句——\emph{这一轮值不值得开通信协商？}——用\textbf{冲突门控}（conflict gating）决定 $g^t\in\{0,1\}$：$g^t{=}0$ 跳过协商、只刷新共识参考；$g^t{=}1$ 才进入 MACPO 原有协商分支。实现体叫 \textbf{RL-MACPO}（在 MACPO 代码栈上加门控；RL 只负责调节惩罚系数，\textbf{不是}主要创新点）。

\subsection{四个研究问题（后文实验按此组织）}
\begin{enumerate}
  \item \textbf{Q1 优化质量}：相同评估预算下，门控版能否不比原始 MACPO 差、甚至更好？
  \item \textbf{Q2 通信削减}：协商触发率能否从 100\% 降到一小部分？
  \item \textbf{Q3 机制归因}：增益来自门控还是 RL？门控三层里哪一层必要？
  \item \textbf{Q4 工程迁移}：经济调度、IEEE 电网等场景是否仍有效？
\end{enumerate}

\subsection{MACPO 外循环（读懂后文公式的前提）}
每个外循环 $t$ 大致为：
\begin{enumerate}[label=\arabic*.]
  \item 各智能体用局部搜索（LLSO/CSO）优化\textbf{带惩罚的局部目标} $h_i$；
  \item 计算冲突指标 CI，经\textbf{门控}得 $g^t$；
  \item 若 $g^t{=}1$：与邻居\textbf{协商}共享变量；若 $g^t{=}0$：\textbf{跳过协商}，仅更新共识参考 $r_d$；
  \item 重复直至评估次数（eva）用尽。
\end{enumerate}
\textbf{评估预算}：两种方法用相同随机种子配对、相同 eva 上限，保证对比公平。

%==============================================================================
\section{术语表：代码名与中文含义}\label{sec:glossary}
%==============================================================================
论文与日志里常出现英文配置名，下表是\textbf{第一次出现时的标准中文叫法}；后文实验优先用中文，括号内保留代码名便于查日志。

\begin{center}\small
\begin{longtable}{@{}p{3.2cm}p{4.2cm}p{6.8cm}@{}}
\toprule
\textbf{代码/论文名} & \textbf{中文名称} & \textbf{含义} \\
\midrule
\endhead
MACPO / always-on &
\textbf{原始 MACPO（每轮必通信）} &
门控恒开 $g^t\equiv 1$，每外循环都协商；\emph{always-on} 即“通信常开”，不是算法正式名称。 \\
\addlinespace
RL-MACPO / Full &
\textbf{完整版 RL-MACPO} &
门控 + RL 惩罚调节 + 变量筛选（0.9/0.7/0.5）；论文主表默认配置；\emph{Full} 指“全开栈”，非 “完整实验” 的泛称。 \\
\addlinespace
RL Only &
\textbf{仅 RL 惩罚（无门控）} &
保留 PPO 调 $\alpha_i$，但 $g^t\equiv 1$，\textbf{没有}冲突门控；用于证明“单靠 RL 不够”。 \\
\addlinespace
NoGating &
\textbf{无门控} &
消融：去掉门控，其余与完整版相同。 \\
\addlinespace
V0 / V0\_AlwaysOn &
\textbf{门控变体 0：每轮必通信} &
门控消融基线，等价于原始 MACPO 通信策略。 \\
V1 & \textbf{变体 1：固定绝对阈值} & 仅当 CI 超过常数 $\theta$ 才通信；无 fail-safe。 \\
V2 & \textbf{变体 2：相对阈值} & CI 高于本阶段 EMA 基线 $\lambda$ 倍才通信；无 fail-safe。 \\
V3 & \textbf{变体 3：相对阈值 + fail-safe} & 本文门控核心；静默超过 $K$ 轮强制通信一次。 \\
\addlinespace
Layer1 / Layer2 / Layer3 &
\textbf{门控第 1/2/3 层} &
Layer1=最小间隔 $\Delta_{\min}$；Layer2=相对阈值（+fail-safe）；Layer3=相位随机采样。 \\
\addlinespace
Periodic-$K$ &
\textbf{固定周期通信} &
不看冲突，每 $K$ 轮强制通信一次（无反馈基线）。 \\
\addlinespace
$f_{\mathrm{pure}}$ / $F$ &
\textbf{纯目标 / 全局目标} &
$f_{\mathrm{pure}}=\sum_i f_i$（\textbf{不含}惩罚，主对比指标，越小越好）；
$F$ 同义或含惩罚的归档 endpoint，读表时以 caption 为准。 \\
\addlinespace
eva / eva\_count &
\textbf{评估次数} &
黑盒目标函数被调用的总次数；与“进化代数”不同。 \\
\addlinespace
$\bar p_{\mathrm{comm}}$ &
\textbf{通信触发率} &
外循环中 $g^t{=}1$ 的比例；原始 MACPO 为 100\%。 \\
\addlinespace
F1--F18 &
\textbf{MACPO NDO 基准函数族} &
F1--F6：20 智能体链式图；F7--F18：更大规模链/网格；来自 MACPO 原文基准套件。 \\
\bottomrule
\end{longtable}
\end{center}

\subsection{变量速查（读伪代码与公式时用）}\label{sec:varindex}
下表列出文中\textbf{高频符号与日志字段}。伪代码块内仅写英文变量名；中文含义见本表或各算法后的「变量说明」。
\par\smallskip\noindent\emph{符号约定：}计数一律写 $n_{\mathrm{win}}$、$n_{\mathrm{comm}}$ 等，或直接用中文描述；\textbf{不使用}数学中的 `\#` 集合计数写法。日志里的 \texttt{COST\_STATS} 是\textbf{字面字段名}，不是 C/ Python 的注释符。
\par\smallskip

\begin{center}\small
\begin{longtable}{@{}p{3.0cm}p{4.0cm}p{7.2cm}@{}}
\toprule
\textbf{符号/字段} & \textbf{英文全称或代码名} & \textbf{含义} \\
\midrule
\endhead
$t$ & outer-loop index & 外循环序号 \\
$g^t$, $g_i^t$ & gate & 第 $t$ 轮是否开通信：1=协商，0=跳过 \\
$g_{\mathrm{global}}^t$ & global gate & 各智能体本地门的最大值（式\eqref{eq:gate}） \\
$T$ & total outer loops & 一次 run 的外循环总数 \\
$n_{\mathrm{comm}}$ & comm counter & $g_{\mathrm{global}}^t{=}1$ 的外循环个数 \\
$\bar p_{\mathrm{comm}}$ & mean trigger rate & $n_{\mathrm{comm}}/T$，通信触发率 \\
$f_{\mathrm{pure}}$ & pure objective & $\sum_i f_i$，\textbf{不含}惩罚，主对比指标 \\
$F$ & global / penalized objective & 全局目标；归档日志有时含惩罚项 \\
$h_i$ & penalized local objective & $f_i+\alpha_i c_i$，局部搜索实际最小化目标 \\
$\alpha_i$ & penalty coefficient & 惩罚强度 \\
$c_i$, CI & conflict index & 共享变量不一致程度；门控输入 \\
$r_d$ & consensus reference & 共享维 $d$ 的共识参考值 \\
$x_d$ & local best on dim $d$ & 智能体在维 $d$ 上的局部最优 \\
$R_d$ & search range & 维 $d$ 的搜索范围（归一化用） \\
$\lambda$ & rel.\ threshold multiplier & 相对阈值系数（默认 1.2） \\
$K$ & fail-safe period & 最长静默轮数，超过则强制通信 \\
$\Delta_{\min}$ & min interval & 两次通信间最少间隔外循环数 \\
$\beta,\gamma$ & EMA coefficients & CI 平滑系数（0.7）与阶段基线 EMA（0.9） \\
\addlinespace
\texttt{run\_id} & run index & 第几次独立重复实验（1--25） \\
\texttt{MACPO\_PAIR\_SEED} & paired seed & 环境变量：配对随机种子，保证 MACPO 与 RL-MACPO 可比 \\
\texttt{eva\_count} & evaluation count & 黑盒目标函数被调用的总次数（评估预算） \\
\texttt{comm\_rate} & communication rate & 日志中的通信触发率，同 $\bar p_{\mathrm{comm}}$ \\
\texttt{wall\_clock\_s} & wall-clock seconds & 单次 run 墙钟时间（秒） \\
\texttt{COST\_STATS} & log footer tag & 轨迹文件\textbf{末尾}的汇总行标记（字面字符串，非注释符） \\
\addlinespace
w/t/l & win / tie / loss & 显著性检验下 RL-MACPO 相对 MACPO 的胜/平/负函数个数 \\
$\Delta_k$ & paired improvement & 第 $k$ 次 run 上 MACPO 与 RL-MACPO 的 $f_{\mathrm{pure}}$ 相对差 \\
FES$(e)$ & fitness vs evaluations & 累计 eva $\le e$ 时所见最优 $f_{\mathrm{pure}}$ \\
\bottomrule
\end{longtable}
\end{center}

%==============================================================================
\section{核心公式：含义与论文出处}\label{sec:notation}
%==============================================================================
下列公式在多篇实验中反复出现。\textbf{左列编号}供后文引用；\textbf{出处}对应 \texttt{conference\_new\_ready.tex} 章节。

\paragraph{(1) 全局 NDO 问题}\label{eq:global-source}
\textit{出处：论文 Section~III（NDO Formulation），式 (1)--(2)。}
\begin{equation}
  \min_{\mathbf{X}} F(\mathbf{X}) = \sum_{i=1}^{n} f_i(\mathbf{X}_i),
  \qquad X_i = \mathrm{col}\{x_d\}_{d\in\Theta_i}
  \label{eq:global}
\end{equation}
$n$ 个智能体各优化局部块 $X_i$；$\Theta_i$ 含私有变量与和邻居共享的边变量。约束为共享维度上 $x_i^d=x_j^d$（MACPO 用惩罚+协商近似满足）。

\paragraph{(2) MACPO 局部惩罚目标}\label{eq:penalized-source}
\textit{出处：论文 Section~V（MACPO Instantiation）；MACPO 原文。}
\begin{equation}
  h_i(\mathbf{X}_i) = f_i(\mathbf{X}_i) + \alpha_i\, c_i(\mathbf{X}_i)
  \label{eq:penalized}
\end{equation}
局部搜索实际最小化 $h_i$：$f_i$ 是真实子目标，$c_i$ 度量共享变量与共识参考 $r_d$ 的不一致；$\alpha_i$ 为惩罚强度（可固定或由 RL/EMA 调节）。

\paragraph{(3) 冲突指标 CI（门控输入）}\label{eq:ci-source}
\textit{出处：论文 Section~IV 式 (4)--(5)；概念来自 MACPO 的 conflict index，本文改为仅用位置差（黑盒可用）。}
\begin{align}
  c_d^t &= \beta c_d^{t-1} + (1-\beta)\frac{|x_d^t - r_d^t|}{R_d},
  \label{eq:ci_dim}\\
  \mathrm{CI}_i^t &= \frac{1}{|\mathcal{D}_i|}\sum_{d\in\mathcal{D}_i} c_d^t
  \label{eq:ci}
\end{align}
$x_d^t$：智能体局部最优；$r_d^t$：共识参考；$R_d$：搜索范围；$\mathcal{D}_i$：智能体 $i$ 上活跃的重叠维。CI 越大表示共享变量越不一致。

\paragraph{(4) 阶段基线（相对阈值用）}\label{eq:ci_ema-source}
\textit{出处：论文 Section~IV，阈值约束 $\mathcal{C}_{\mathrm{thr}}$ 中的 EMA。}
\begin{equation}
  \hat\mu_{\mathrm{CI},i}^t = \gamma \hat\mu_{\mathrm{CI},i}^{t-1} + (1-\gamma)\mathrm{CI}_i^t
  \label{eq:ci_ema}
\end{equation}
$\hat\mu_{\mathrm{CI},i}^t$ 是“本阶段正常冲突水平”；当前 CI 显著高于它时才倾向于开通信。

\paragraph{(5) 冲突门控（本文核心）}\label{eq:gate-source}
\textit{出处：论文 Section~IV 式 (6)--(7)，Algorithm~1。}
\begin{align}
  g_i^t &= \mathbf{1}\bigl[\mathcal{C}_{\mathrm{int}}^t \land \mathcal{C}_{\mathrm{thr}}^t \land \mathcal{C}_{\mathrm{phase}}^t\bigr],
  \label{eq:local_gate}\\
  g_{\mathrm{global}}^t &= \max_i g_i^t
  \label{eq:gate}
\end{align}
三层条件（须\textbf{同时}满足才开本地门）：
\begin{itemize}
  \item $\mathcal{C}_{\mathrm{int}}$：距上次通信 $\ge \Delta_{\min}$（默认 10 轮）；
  \item $\mathcal{C}_{\mathrm{thr}}$：$\mathrm{CI}_i^t > \lambda\hat\mu_{\mathrm{CI},i}^t$ \textbf{或} fail-safe：距上次通信 $\ge K$（默认 $K{=}10$）；
  \item $\mathcal{C}_{\mathrm{phase}}$：以概率 $p_{\mathrm{phase}}$ 抽样通过（早/中/晚：0.6/0.3/0.15）。
\end{itemize}
全局门取各智能体最大值：任一智能体请求通信则本轮全体进入协商。

\paragraph{(6) 通信触发率与削减率}\label{eq:comm-source}
\textit{出处：论文 Section~III（Communication-Timing Problem）、实验 Q2。}
\begin{align}
  \bar p_{\mathrm{comm}} &= \frac{1}{T}\sum_{t=1}^{T} g_{\mathrm{global}}^t,
  \label{eq:pcomm}\\
  \text{通信削减率} &= \bigl(1-\bar p_{\mathrm{comm}}^{\mathrm{RL-MACPO}}\bigr)\times 100\%
  \quad\text{（相对原始 MACPO 的 100\%）}
  \label{eq:comm}
\end{align}

\paragraph{主实验默认参数}
\textit{出处：论文 Table~II（gate defaults）。}
$\beta{=}0.7$，$\gamma{=}0.9$，$\lambda{=}1.2$，$\Delta_{\min}{=}10$，$K{=}10$（IEEE 试点 $K{=}2$）；变量筛选比例 0.9/0.7/0.5（按优化阶段取 Top 比例共享维参与协商）。

%==============================================================================
\section{实验索引：26 条实验与论文问题对应}\label{sec:roadmap}
%==============================================================================
\begin{center}\scriptsize
\begin{longtable}{@{}clp{7.5cm}@{}}
\toprule
\textbf{编号} & \textbf{论文问题} & \textbf{一句话在测什么} \\
\midrule
\endhead
1--2 & Q1 & F1--F18 上完整版 vs 原始 MACPO 的最终 $f_{\mathrm{pure}}$ \\
3 & Q1 补充 & 收敛曲线（FES），定性看是否更快/更低 \\
4 & 补充 & 墙钟时间（非主指标） \\
5--6 & Q2 & 通信触发率 + eva 预算是否公平 \\
7--12 & Q2/Q3 & 门控变体、$\lambda$ 敏感、周期基线、CI 分箱、Pareto \\
13, 19--21 & Q3 诊断 & 单 run / 25-run 轨迹：CI、$\alpha$、漂移 \\
14 & 外部参考 & MASOIE 跨范式对比（非公平 matched-budget） \\
15--18 & Q3 & 惩罚控制器、门控层、相位/$\eta$、RL 网络结构 \\
22--25 & Q4 & MAED-13、资源调度、EV、IEEE 30/57/118 \\
26 & 可扩展性 & 智能体数增多时的 pilot（runs 较少，谨慎解读） \\
\bottomrule
\end{longtable}
\end{center}

%==============================================================================
\section{实验 1：F1--F6 优化质量主对比（完整版 RL-MACPO vs 原始 MACPO）}
%==============================================================================
\subsection{实验说明}
\textbf{对应论文 Q1。}在相同 eva 预算下，比较\textbf{完整版 RL-MACPO}（冲突门控 + RL 惩罚 + 变量筛选）与\textbf{原始 MACPO（每轮必通信）}的最终纯目标 $f_{\mathrm{pure}}$。另列\textbf{仅 RL 惩罚（无门控）}、GFPDO$^\dagger$、DPSO$^\dagger$ 作参考（后两者来自 MACPO 论文表格布局，不做显著性检验）。基准：20 智能体链式图，$D{=}905$，函数 F1--F6。

\subsection{实验计算公式}
\begin{itemize}
  \item 优化问题：式\eqref{eq:global}（Section~\ref{sec:notation}）；
  \item 门控逻辑：式\eqref{eq:local_gate}--\eqref{eq:gate}；\textbf{原始 MACPO} 等价于恒令 $g^t\equiv 1$；
  \item 门控开启（$g^t{=}1$）时走 MACPO 协商；关闭时只刷新 $r_d$，继续局部 LLSO/CSO。
\end{itemize}

\subsection{算法伪代码}
\begin{algorithm}[H]
\caption{F1--F6 main comparison (paired runs)}
\begin{algorithmic}[1]
\For{$\textit{run\_id} = 1$ \textbf{to} $25$}
  \State $\textsc{PairSeed} \gets \textit{run\_id}$ \Comment{env \texttt{MACPO\_PAIR\_SEED}}
  \State Run baseline MACPO with $g^t \equiv 1$ until \textit{eva\_limit}
  \State Run full RL-MACPO until the same \textit{eva\_limit}
  \State Log final $f_{\mathrm{pure}}$, $\bar p_{\mathrm{comm}}$, \textit{eva\_count}
\EndFor
\State LLSO: paired Wilcoxon test; CSO: Mann--Whitney U test ($\alpha=0.05$)
\end{algorithmic}
\end{algorithm}
\varnote{
\texttt{run\_id}：第几次独立 run（1--25）。
\textsc{PairSeed} / \texttt{MACPO\_PAIR\_SEED}：配对随机种子。
$g^t\equiv 1$：原始 MACPO，每轮必通信。
\textit{eva\_limit}：评估次数上限（两种方法相同）。
$f_{\mathrm{pure}}$：纯目标；$\bar p_{\mathrm{comm}}$：通信触发率；\textit{eva\_count}：实际调用目标函数的次数。
}

\subsection{结果评估指标公式}
\begin{align}
  f_{\mathrm{pure}} &= \sum_i f_i(\mathbf{X}_i)
  \quad\text{（纯目标 $f_{\mathrm{pure}}$：不含惩罚项，越小越好）} \label{eq:fpure}\\
  \text{w/t/l} &= n_{\mathrm{win}}/n_{\mathrm{tie}}/n_{\mathrm{loss}}
  \quad\text{（胜/平/负：RL-MACPO 相对 MACPO 在 0.05 水平下显著更优/无显著差异/显著更差的函数个数）} \label{eq:wtl}
\end{align}
配对 Wilcoxon 基于差值 $\Delta_k = f_{\mathrm{pure},k}^{\mathrm{MACPO}} - f_{\mathrm{pure},k}^{\mathrm{RL-MACPO}}$（$\Delta_k>0$ 表示 RL-MACPO 更优）。

\subsection{实验设置与参数}
\begin{center}\small
\begin{tabular}{@{}ll@{}}
\toprule
项目 & 值 \\
\midrule
函数 & F1--F6（同质 F1--F3，异质 F4--F6） \\
优化器 & LLSO（主）、CSO（次） \\
Runs & 25 配对独立运行 \\
RL-MACPO 配置 & 完整版：门控（式\eqref{eq:local_gate}）+ PPO 调 $\alpha$ + 变量筛选 0.9/0.7/0.5 \\
\bottomrule
\end{tabular}
\end{center}

\subsection{实验结果}
\begin{center}\scriptsize
\begin{tabular}{@{}lccccc@{}}
\toprule
函数 & MACPO (LLSO) & RL-MACPO (LLSO) & $p$ 值 & w/t/l & 通信削减 \\
\midrule
F1 & $1.98\times10^7$ & $\mathbf{6.43\times10^6}$ & $2.98\times10^{-8}$ & 6/0/0 & 79.2\% \\
F2 & $5.05\times10^4$ & $\mathbf{3.32\times10^3}$ & $2.98\times10^{-8}$ & — & 77.7\% \\
F3 & $4.03\times10^9$ & $\mathbf{6.35\times10^2}$ & $2.98\times10^{-8}$ & — & 91.7\% \\
F4 & $3.17\times10^6$ & $\mathbf{2.58\times10^6}$ & $8.94\times10^{-8}$ & — & 91.7\% \\
F5 & $1.97\times10^8$ & $\mathbf{3.50\times10^6}$ & $2.98\times10^{-8}$ & — & 91.7\% \\
F6 & $1.62\times10^8$ & $\mathbf{2.47\times10^3}$ & $2.98\times10^{-8}$ & — & 91.7\% \\
\bottomrule
\end{tabular}
\end{center}
\textbf{解读}：完整版 RL-MACPO 六函数全胜；\textbf{仅 RL 惩罚（无门控）}在 F3/F5 弱于完整版，说明增益主要来自\textbf{冲突门控}，而非单靠 RL 调惩罚。

\tryfig{media/F1_F6_panel.pdf}{实验 1 配套图：F1--F6 单 run 配对 best-so-far 收敛曲线}

%==============================================================================
\section{实验 2：F7--F18 优化质量主对比}
%==============================================================================
\subsection{实验说明}
\textbf{对应论文 Q1 扩展。}将实验 1 推广至 F7--F18（40/60 智能体等更大 NDO 实例），检验门控在大规模上是否仍有效。

\subsection{实验计算公式}
同 Section~\ref{sec:notation}；注意 MACPO 归档 endpoint 有时记为带惩罚的 $F$，RL-MACPO 主表报 $f_{\mathrm{pure}}$（纯目标，式\eqref{eq:global} 右端不含 $\alpha c$）。

\subsection{算法伪代码}
同实验 1，函数集合改为 F7--F18，其余不变。

\subsection{结果评估指标公式}
\begin{equation}
  \text{报告格式：}\;\overline{F \pm \sigma},\quad \text{w/t/l 按 12 对函数统计}
\end{equation}

\subsection{实验设置与参数}
25 配对 runs；数据源 MACPO\_original（MACPO）+ MACPO2\_deployment（RL）；LLSO/CSO 双优化器。

\subsection{实验结果}
LLSO 与 CSO 均为 \textbf{12/0/0} 全胜。典型：F9 MACPO $3.05\times10^{11}$ vs RL $1.01\times10^3$；F15 MACPO $1.24\times10^{12}$ vs RL $8.41\times10^2$。高冲突 Rosenbrock 相关实例差距最大。

%==============================================================================
\section{实验 3：F1--F6 收敛曲线（FES 面板）}
%==============================================================================
\subsection{实验说明}
从 25-run 批次中各取\textbf{一组配对代表 run}，绘制 best-so-far 纯目标随累计评估次数的收敛轨迹，定性展示 Q1 结论。

\subsection{实验计算公式}
\begin{equation}
  \mathrm{FES}(e) = \min\bigl\{ f_{\mathrm{pure}} \mid \text{累计 eva} \le e \bigr\}
\end{equation}
其中 $e$ 为累计评估次数，$f_{\mathrm{pure}}$ 为纯目标（见式\eqref{eq:fpure}）。

\subsection{算法伪代码}
\begin{algorithm}[H]
\caption{Extract FES convergence curves}
\begin{algorithmic}[1]
\For{each benchmark function $F_k \in \{\mathrm{F1},\ldots,\mathrm{F6}\}$}
  \State Select one \textit{run\_id} whose MACPO and RL-MACPO logs are both complete
  \State Build $\mathrm{FES}(e)$ from the \textit{eva} sequence
  \State Interpolate onto a common \textit{eva} grid and plot panel $k$
\EndFor
\end{algorithmic}
\end{algorithm}
\varnote{
$F_k$：第 $k$ 个基准函数（F1--F6）。
$\mathrm{FES}(e)$：累计评估次数不超过 $e$ 时所见最优 $f_{\mathrm{pure}}$（见下式）。
\textit{eva}：每次黑盒目标评估对应的累计计数。
}

\subsection{结果评估指标公式}
纵轴 $\mathrm{FES}$（越低越好）；横轴累计 eva；不报告 $p$ 值（定性图，定量见实验 1）。

\subsection{实验设置与参数}
出图脚本 \texttt{plot\_fes\_f1\_f6\_panel.py}；日志 MACPO\_original + RL-output\_runs25。

\subsection{实验结果}
\tryfig{media/F1_F6_panel.pdf}{F1--F6 配对代表 run 的 FES 收敛面板}

%==============================================================================
\section{实验 4：墙钟时间对比（F1--F18）}
%==============================================================================
\subsection{实验说明}
在开发机上测量单次 run 墙钟时间，\textbf{非主终点}，用于说明 RL 推理与门控 bookkeeping 的 indicative 开销。

\subsection{实验计算公式}
\begin{equation}
  \text{Overhead ratio} \approx T_{\mathrm{RL\mbox{-}MACPO}} / T_{\mathrm{MACPO}}
\end{equation}

\subsection{算法伪代码}
同实验 1 运行流程；在相同 25-run 批次上额外记录 \texttt{wall\_clock\_s}（硬编码于 tex，非 fitness 日志统一字段）。

\subsection{结果评估指标公式}
报告各函数 mean wall-clock（秒）；F1 LLSO 例：MACPO 14.43\,s，RL-MACPO 26.95\,s（约 1.9$\times$）。

\subsection{实验设置与参数}
开发机 indicative 测量；与 Table \texttt{tab:wall\_time\_f1\_f18} 一致。

\subsection{实验结果}
F1--F6 上 RL-MACPO 约为 MACPO 的 1.7--2.1 倍；F7--F18 比值类似。通信削减带来的协商次数下降\textbf{未完全抵消} RL 模块开销，但 fitness 主指标仍全胜。

%==============================================================================
\section{实验 5：F1--F18 通信触发率统计}
%==============================================================================
\subsection{实验说明}
\textbf{对应论文 Q2。}量化\textbf{完整版 RL-MACPO}相对\textbf{原始 MACPO（每轮必通信）}的外循环协商触发比例 $\bar p_{\mathrm{comm}}$（式\eqref{eq:pcomm}--\eqref{eq:comm}），验证通信能否大幅削减且 Q1 质量不降。

\subsection{实验计算公式}
式\eqref{eq:pcomm}--\eqref{eq:comm}；\textbf{原始 MACPO} 基线 $\bar p_{\mathrm{comm}} = 1$（100\%）。

\subsection{算法伪代码}
\begin{algorithm}[H]
\caption{Communication trigger rate from a trajectory}
\begin{algorithmic}[1]
\State $n_{\mathrm{comm}} \gets 0$
\For{each outer loop $t = 1$ \textbf{to} $T$}
  \If{$g_{\mathrm{global}}^t = 1$}
    \State $n_{\mathrm{comm}} \gets n_{\mathrm{comm}} + 1$
  \EndIf
\EndFor
\State $\bar p_{\mathrm{comm}} \gets n_{\mathrm{comm}} / T$
\end{algorithmic}
\end{algorithm}
\varnote{
$T$：外循环总数。
$g_{\mathrm{global}}^t$：第 $t$ 轮全局门控（式\eqref{eq:gate}）。
$n_{\mathrm{comm}}$：触发协商的外循环个数。
$\bar p_{\mathrm{comm}}$：通信触发率（式\eqref{eq:pcomm}）。
}

\subsection{结果评估指标公式}
\begin{align}
  \text{削减率} &= (1 - \bar p_{\mathrm{comm}}^{\mathrm{RL}})\times 100\% \\
  \text{eva 一致性} &: \;|\mathrm{eva}_{\mathrm{RL}} - \mathrm{eva}_{\mathrm{MACPO}}| / \mathrm{eva}_{\mathrm{MACPO}} \approx 3.2\%\;\text{（F1--F6）}
\end{align}

\subsection{实验设置与参数}
LLSO，25 配对 runs；\textbf{完整版 RL-MACPO}；脚本 \texttt{run\_comm\_rate\_f1\_f18.sh}。

\subsection{实验结果}
\begin{center}\small
\begin{tabular}{@{}lc@{}}
\toprule
函数组 & RL-MACPO 触发率 \\
\midrule
F1 & 20.8\% \\
F2 & 22.3$\pm$4.1\% \\
F3--F6 & 8.3\%（削减 91.7\%） \\
F7--F12, F14--F18 & 8.0\% \\
F13 & 7.7$\pm$1.6\% \\
\bottomrule
\end{tabular}
\end{center}
F3--F18 约每 12 个外循环协商 1 次（fail-safe + 间隔 + 相位采样共同作用）。

%==============================================================================
\section{实验 6：F1--F6 通信率与 eva 预算一致性}
%==============================================================================
\subsection{实验说明}
同步报告通信率与评估次数，排除``通信少是因为提前终止''的质疑。

\subsection{实验计算公式}
\begin{equation}
  \textit{eva\_count} = \text{一次 run 中目标函数被调用的总次数},\quad
  \Delta_{\mathrm{eva}} = \frac{\textit{eva\_count}_{\mathrm{RL}} - \textit{eva\_count}_{\mathrm{MACPO}}}{\textit{eva\_count}_{\mathrm{MACPO}}}
\end{equation}
$\textit{eva\_count}$ 即评估预算实际消耗；$\Delta_{\mathrm{eva}}$ 衡量两种方法评估次数相对差异。

\subsection{算法伪代码}
从轨迹文件\textbf{末尾}读取以 \texttt{COST\_STATS} 开头的汇总行（日志格式标记，不是代码注释），提取 \textit{eva\_count} 与 \textit{comm\_rate}。
\varnote{
\texttt{COST\_STATS}：程序写入日志末尾的统计行前缀；后接 \textit{eva\_count}、\textit{comm\_rate} 等字段。
\textit{comm\_rate}：与 $\bar p_{\mathrm{comm}}$ 同义。
}

\subsection{结果评估指标公式}
F1--F6 pooled：MACPO eva $=145729$，RL eva $=150348$，$\Delta_{\mathrm{eva}}\approx +3.2\%$。

\subsection{实验设置与参数}
与实验 5 同一批\textbf{完整版}日志（Exp4 Selection\_0.9\_0.7\_0.5）。

\subsection{实验结果}
通信削减 77--92\% 的同时，RL 评估次数略\textbf{高于} MACPO（因门控跳过协商后局部搜索路径不同），差异在 3\% 量级，不构成预算不公平。

%==============================================================================
\section{实验 7：门控变体对比（V0--V3）}
%==============================================================================
\subsection{实验说明}
\textbf{对应论文 Q3（门控设计）。}对比四种门控变体（见术语表），证明\textbf{相对阈值 + fail-safe（变体 3）}的必要性：若只用固定绝对阈值（变体 1），通信率可降为 0，共识参考冻结，优化失败。

\subsection{实验计算公式}
在式\eqref{eq:local_gate} 的 $\mathcal{C}_{\mathrm{thr}}$ 上换规则：
\begin{align}
  \mathcal{C}_{\mathrm{thr}}^{\mathrm{rel}} &: \mathrm{CI}_i^t > \lambda \hat\mu_{\mathrm{CI},i}^t \;\text{或}\; t - t_{\mathrm{last\_comm}} \ge K \label{eq:thr_rel}\\
  \mathcal{C}_{\mathrm{thr}}^{\mathrm{abs}} &: \mathrm{CI}_i^t > \theta_{\mathrm{fixed}} \quad\text{（无常数 fail-safe）} \label{eq:thr_abs}
\end{align}
式\eqref{eq:thr_rel} 即论文默认；式\eqref{eq:thr_abs} 为变体 1。

\subsection{算法伪代码}
\begin{algorithm}[H]
\caption{Gate variants V0--V3}
\begin{algorithmic}[1]
\State V0: enforce $g^t \equiv 1$ \Comment{baseline MACPO, always communicate}
\State V1: $\mathcal{C}_{\mathrm{thr}}^{\mathrm{abs}}$ only, no fail-safe \Comment{Eq.~\eqref{eq:thr_abs}}
\State V2: relative threshold only, no fail-safe \Comment{first clause of Eq.~\eqref{eq:thr_rel}}
\State V3: relative threshold + fail-safe $K$ \Comment{full gate, Eq.~\eqref{eq:thr_rel}}
\end{algorithmic}
\end{algorithm}
\varnote{
V0--V3：门控消融变体（见术语表）。
$g^t$：外循环通信门。
$\mathcal{C}_{\mathrm{thr}}^{\mathrm{abs}}$：固定绝对阈值条件；$\mathcal{C}_{\mathrm{thr}}^{\mathrm{rel}}$：相对阈值或 fail-safe。
$K$：fail-safe 周期（默认 10）。
}

\subsection{结果评估指标公式}
通信率均值$\pm$标准差；零通信 run 数；协商维度数；墙钟（秒）。

\subsection{实验设置与参数}
F1/F2/F5 混合统计；30 runs；LLSO。

\subsection{实验结果}
\begin{center}\small
\begin{tabular}{@{}lcccc@{}}
\toprule
变体 & 零通信 run & 通信率 & 墙钟(s) & 削减 \\
\midrule
V0 每轮必通信 & 0/30 & 100\% & 6.79 & 0\% \\
V1 固定绝对阈值 & 30/30 & 0\% & 6.41 & 100\% \\
V2 相对阈值 & 0/30 & 50.4\% & 6.43 & 50\% \\
V3 相对+fail-safe & 0/30 & 46.0\% & 6.61 & 54\% \\
\bottomrule
\end{tabular}
\end{center}
V1 完全静默导致共识冻结；V3 在削减通信的同时保证非零协商。

%==============================================================================
\section{实验 8：相对阈值 $\lambda$ 敏感性（F3/F5）}
%==============================================================================
\subsection{实验说明}
扫描 $\lambda \in \{0.8,1.0,1.2,1.4\}$，检验默认 $\lambda{=}1.2$ 是否为窄调优最优。

\subsection{实验计算公式}
式\eqref{eq:local_gate} 中 $\mathcal{C}_{\mathrm{thr}}$ 的 $\lambda$ 为扫描变量；其余与\textbf{完整版默认参数}相同，$K{=}10$。

\subsection{算法伪代码}
\begin{algorithm}[H]
\caption{$\lambda$ sensitivity sweep on F3/F5}
\begin{algorithmic}[1]
\For{$\lambda \in \{0.8, 1.0, 1.2, 1.4\}$}
  \For{$\textit{run\_id} = 1$ \textbf{to} $10$}
    \State Fix \textsc{PairSeed}; run F3 or F5 with threshold multiplier $\lambda$
    \State Record $\bar p_{\mathrm{comm}}$ and final $f_{\mathrm{pure}}$
  \EndFor
\EndFor
\end{algorithmic}
\end{algorithm}
\varnote{
$\lambda$：相对阈值系数（式\eqref{eq:local_gate} 中 $\mathcal{C}_{\mathrm{thr}}$）。
\textsc{PairSeed}：配对种子 \texttt{MACPO\_PAIR\_SEED}。
}

\subsection{结果评估指标公式}
主指标：$\bar p_{\mathrm{comm}}$（均值$\pm$标准差）；辅指标：final $F$。

\subsection{实验设置与参数}
F3/F5，LLSO，$n{=}10$/设置，Selection\_0.9\_0.7\_0.5。

\subsection{实验结果}
\begin{center}\small
\begin{tabular}{@{}lcc@{}}
\toprule
$\lambda$ & F3 通信率 & F5 通信率 \\
\midrule
0.8 & 25.0\% & 27.5$\pm$5.0\% \\
1.0 & 8.3\% & 30.8$\pm$5.0\% \\
1.2（默认） & 8.3\% & 15.8$\pm$10.0\% \\
1.4 & 8.3\% & 8.3\% \\
\bottomrule
\end{tabular}
\end{center}
F5 触发率随 $\lambda$ 单调下降；F3 在 $\lambda\ge 1.0$ 时触达 fail-safe 下限 8.3\%。

%==============================================================================
\section{实验 9：Periodic-$K$ 固定周期通信基线}
%==============================================================================
\subsection{实验说明}
与冲突门控对比\textbf{无反馈的固定周期调度}：每 $K$ 个外循环强制通信一次。

\subsection{实验计算公式}
\begin{equation}
  g^t = \mathbf{1}\bigl[(t \bmod K) = 0\bigr]
\end{equation}

\subsection{算法伪代码}
\begin{algorithm}[H]
\caption{Periodic-$K$ communication baseline}
\begin{algorithmic}[1]
\For{each outer loop $t$}
  \If{$t - t_{\mathrm{last\_comm}} \ge K$}
    \State $g^t \gets 1$; run MACPO negotiation
  \Else
    \State $g^t \gets 0$; skip negotiation
  \EndIf
\EndFor
\end{algorithmic}
\end{algorithm}
\varnote{
$t_{\mathrm{last\_comm}}$：上次触发通信的外循环序号。
$K$：固定周期间隔；理想通信率约为 $1/K$。
$g^t$：本轮是否进入协商分支。
}

\subsection{结果评估指标公式}
通信率 $= 1/K$（理想周期）；final $f_{\mathrm{pure}}$ 均值$\pm$标准差；与\textbf{完整版 RL-MACPO} 比较 Pareto 位置。

\subsection{实验设置与参数}
F1/F2/F5；PeriodicK2/K3/K5；pilot 10-run（F125）或 25-run（F1--F6 扩展）。

\subsection{实验结果}
\tryfig{media/comm_fitness_pareto_f125.pdf}{实验 9/12：通信率--适应度 Pareto 前沿（F1/F2/F5）}
\textbf{完整版 RL-MACPO}在相近通信预算下 fitness 更低，优于固定周期策略。

%==============================================================================
\section{实验 10：阈值触发通信基线（F1--F6）}
%==============================================================================
\subsection{实验说明}
对比\textbf{完整版}、Periodic-$K$、固定绝对阈值、无 fail-safe、相对阈值+fail-safe 等通信调度策略（均替换式\eqref{eq:local_gate} 中的门控模块，协商与局部搜索不变）。

\subsection{实验计算公式}
见实验 7--9；额外含 FixedThresholdNoFailSafe（$\bar p_{\mathrm{comm}}$ 可为 0）。

\subsection{算法伪代码}
在 MACPO 外循环入口替换门控模块，其余局部搜索与协商不变。

\subsection{结果评估指标公式}
双指标：$\bar p_{\mathrm{comm}}$ 与 final $f_{\mathrm{pure}}$；25 runs/方法/函数。

\subsection{实验设置与参数}
数据源 \texttt{comm\_baselines\_f1\_f6.json}；LLSO。

\subsection{实验结果}
无 fail-safe 时通信率可为 0、fitness 方差极大；\textbf{完整版}在通信--质量权衡上最优。例：F3 完整版通信 8.3\%，FixedThresholdNoFailSafe 通信 0\% 但结果不稳定。

%==============================================================================
\section{实验 11：CI 分箱触发概率统计}
%==============================================================================
\subsection{实验说明}
检验\textbf{低冲突跳过准则}：冲突代理 CI 越低，触发概率应越低。

\subsection{实验计算公式}
\begin{equation}
  P(\text{trigger} \mid \mathrm{CI} \in \mathrm{bin}_q)
  = \frac{n_{\mathrm{trigger},q}}{n_{\mathrm{total},q}}
\end{equation}
其中 $n_{\mathrm{trigger},q}$ 为 CI 落在第 $q$ 箱且 $g^t{=}1$ 的外循环个数，$n_{\mathrm{total},q}$ 为 CI 落在该箱的外循环总数；将 CI 分为五等分位 $\mathrm{bin}_1,\ldots,\mathrm{bin}_5$。

\subsection{算法伪代码}
\begin{algorithm}[H]
\caption{CI-bin conditional trigger probability}
\begin{algorithmic}[1]
\State Load $(\mathrm{CI}^t, g^t)$ sequences from 25 runs
\State Split CI into 5 quantile bins $\mathrm{bin}_1,\ldots,\mathrm{bin}_5$
\For{each bin $q$}
  \State Compute $P(\text{trigger} \mid \mathrm{CI} \in \mathrm{bin}_q)$ and plot
\EndFor
\end{algorithmic}
\end{algorithm}
\varnote{
$\mathrm{CI}^t$：第 $t$ 轮冲突指标（式\eqref{eq:ci}）。
$g^t$：第 $t$ 轮门控；trigger 指 $g^t{=}1$。
$\mathrm{bin}_q$：按 CI 五等分位的第 $q$ 箱。
}

\subsection{结果评估指标公式}
条件触发概率 $P(\text{trigger}\mid\mathrm{bin}_q)$；期望：$\mathrm{bin}_1 > \mathrm{bin}_2 \approx \mathrm{bin}_3$（低冲突箱触发更少）。

\subsection{实验设置与参数}
F1--F6，25-run，日志 \texttt{RL-output\_runs25\_rho/}；脚本 \texttt{plot\_ci\_bin\_trigger.py}。

\subsection{实验结果}
\tryfig{media/ci_bin_trigger_F1_F6.pdf}{F1--F6 CI 五等分位条件触发概率}

%==============================================================================
\section{实验 12：通信--适应度 Pareto 前沿（F1/F2/F5）}
%==============================================================================
\subsection{实验说明}
在\textbf{预算形式}目标下可视化不同通信策略的 Pareto 位置：
$\min F(\mathbf{X}^T)$ s.t.\ $\bar p_{\mathrm{comm}} \le B$。

\subsection{实验计算公式}
横轴 $\bar p_{\mathrm{comm}}$，纵轴 terminal $f_{\mathrm{pure}}$；各方法为 $(\bar p, \overline{f})$ 散点。

\subsection{算法伪代码}
汇总\textbf{完整版}、PeriodicK2/3/5、原始 MACPO 等方法的 $\bar p_{\mathrm{comm}}$ 与 $f_{\mathrm{pure}}$，绘制 Pareto 图。

\subsection{结果评估指标公式}
越靠近\textbf{左下角}（低通信、低 fitness）越优。

\subsection{实验设置与参数}
F1/F2/F5 pilot；脚本 \texttt{plot\_comm\_fitness\_pareto.py}。

\subsection{实验结果}
\tryfig{media/comm_fitness_pareto_f125.pdf}{通信率 vs 最终适应度 Pareto 图}
\textbf{完整版 RL-MACPO（冲突门控）}占据更优前沿区域。

%==============================================================================
\section{实验 13：F3/F5 微观诊断（冲突--惩罚耦合）}
%==============================================================================
\subsection{实验说明}
单 run 尺度展示冲突代理、惩罚系数 $\alpha$、比率 $\rho$ 的\textbf{时序耦合}，解释 RL 控制器行为（非优越性证明）。

\subsection{实验计算公式}
\begin{equation}
  h_i = f_i + \alpha_i c_i,\quad \rho_i = \alpha_i / \alpha_i^{\mathrm{base}}
\end{equation}

\subsection{算法伪代码}
从单条 F3 或 F5 轨迹日志提取 $(e, \alpha, \rho, \mathrm{CI})$ 随 eva 序列并叠加绘图。

\subsection{结果评估指标公式}
定性时序图；无假设检验。

\subsection{实验设置与参数}
单 run；脚本 \texttt{plot\_f3\_f5\_conflict\_penalty\_micro.py}。

\subsection{实验结果}
\tryfig{media/f3_f5_conflict_alpha_micro.pdf}{F3/F5 单 run 冲突--$\alpha$ 微观耦合}

%==============================================================================
\section{实验 14：MASOIE 跨范式外部参考（F3/F5）}
%==============================================================================
\subsection{实验说明}
引入 MASOIE（共识型 DBO）作\textbf{跨范式参考}；拓扑与 MACPO 不匹配，非 matched-budget 公平对比。

\subsection{实验计算公式}
同式\eqref{eq:global}；MASOIE 使用作者可执行文件，无 MACPO 式惩罚协商。

\subsection{算法伪代码}
分别运行\textbf{完整版 RL-MACPO}（25-run）与 MASOIE 官方流程（25-run），比较 final $F$。

\subsection{结果评估指标公式}
final $F$ 均值$\pm$标准差；越低越好；\textbf{不做} Wilcoxon（不同算法栈）。

\subsection{实验设置与参数}
F3/F5，LLSO；MAES-CCSA 开源不可用未纳入。

\subsection{实验结果}
\begin{center}\small
\begin{tabular}{@{}lcc@{}}
\toprule
方法 & F3 final $F$ & F5 final $F$ \\
\midrule
RL-MACPO & $\mathbf{6.35\times10^2 \pm 3.12\times10^1}$ & $\mathbf{3.50\times10^6 \pm 1.40\times10^5}$ \\
MASOIE & $1.04\times10^4 \pm 3.67\times10^3$ & $2.02\times10^7 \pm 1.52\times10^6$ \\
\bottomrule
\end{tabular}
\end{center}

%==============================================================================
\section{实验 15：惩罚控制器对比（F3/F5）}
%==============================================================================
\subsection{实验说明}
\textbf{对应论文 Q3（惩罚控制器）。}在\textbf{门控与变量筛选均固定}为完整版配置的前提下，比较：原始 MACPO 静态惩罚 vs 门控下的 Fixed / EMA / RL 三种 $\alpha_i$ 调节方式。目的：RL 是否必需？——答案应是否定的，门控才是主因。

\subsection{实验计算公式}
在式\eqref{eq:penalized} 中换 $\alpha_i$ 的来源：
\begin{equation}
  h_i = f_i + \alpha_i(\mathrm{state})\, c_i,\quad
  \alpha_i \in \{\alpha^{\mathrm{fixed}}, \alpha^{\mathrm{EMA}}, \alpha^{\mathrm{RL(PPO)}}\}
\end{equation}
门控仍用式\eqref{eq:local_gate}--\eqref{eq:gate}，与实验 15 无关的模块保持不变。

\subsection{算法伪代码}
\begin{algorithm}[H]
\caption{Penalty-controller comparison on F3/F5}
\begin{algorithmic}[1]
\State Fix full-stack gate ($\lambda$, $K$, $p_{\mathrm{phase}}$, variable selection)
\For{each \textit{controller} $\in \{\mathrm{MACPO}, \mathrm{Fixed}, \mathrm{EMA}, \mathrm{RL}\}$}
  \State Run 25 paired trials on F3 and F5
  \State Log best-so-far $f_{\mathrm{pure}}$ and $\bar p_{\mathrm{comm}}$
\EndFor
\State Pairwise Wilcoxon tests among gated controllers
\end{algorithmic}
\end{algorithm}
\varnote{
\textit{controller}：惩罚系数 $\alpha_i$ 的更新方式（Fixed / EMA / RL-PPO）；MACPO 为静态惩罚且无门控。
$p_{\mathrm{phase}}$：相位采样通过概率。
}

\subsection{结果评估指标公式}
best-so-far $f_{\mathrm{pure}}$ + $\bar p_{\mathrm{comm}}$；gated 组间 Wilcoxon 无显著差异。

\subsection{实验设置与参数}
25 runs/控制器/函数；LLSO。

\subsection{实验结果}
\begin{center}\small
\begin{tabular}{@{}lcc|cc@{}}
\toprule
& \multicolumn{2}{c|}{F3} & \multicolumn{2}{c}{F5} \\
控制器 & $F$ & Comm. & $F$ & Comm. \\
\midrule
MACPO & $4.03\times10^9$ & 100\% & $1.97\times10^8$ & 100\% \\
Fixed & $1.16\times10^3$ & 8.3\% & $6.74\times10^6$ & 17.3\% \\
EMA & $1.14\times10^3$ & 8.3\% & $6.94\times10^6$ & 17.3\% \\
RL & $1.16\times10^3$ & 8.3\% & $6.86\times10^6$ & 15.8\% \\
\bottomrule
\end{tabular}
\end{center}
\textbf{结论}：门控固定后，惩罚控制器选择\textbf{统计等价}；主要增益来自通信架构。

%==============================================================================
\section{实验 16：门控层消融（F1--F6）}
%==============================================================================
\subsection{实验说明}
\textbf{对应论文 Q3（门控层消融）。}逐层叠加门控：无门控 $\to$ 仅 Layer1（最小间隔）$\to$ Layer1+2（+相对阈值）$\to$ 完整版（+Layer3 相位采样）。见术语表 Layer1--3 定义。

\subsection{实验计算公式}
对应式\eqref{eq:local_gate} 中条件的逐步启用：
\begin{align}
  \text{无门控}: &\; g^t \equiv 1 \\
  \text{Layer1}: &\;\mathcal{C}_{\mathrm{int}}^t \\
  \text{Layer1+2}: &\;\mathcal{C}_{\mathrm{int}}^t \land \mathcal{C}_{\mathrm{thr}}^t \\
  \text{完整版}: &\;\mathcal{C}_{\mathrm{int}}^t \land \mathcal{C}_{\mathrm{thr}}^t \land \mathcal{C}_{\mathrm{phase}}^t
\end{align}

\subsection{算法伪代码}
在\textbf{完整版}栈上开关各层，25-run 批量运行 F1--F6。

\subsection{结果评估指标公式}
final $F$ 均值/标准差/变异系数 CV；Layer2 缺失时 CV 最高 17.6\%（F1）。

\subsection{实验设置与参数}
25 runs；数据源 \texttt{补做实验\_25runs/}。

\subsection{实验结果}
\tryfig{media/fig_gating_ablation.pdf}{门控层消融：F1--F6 各配置最终 fitness}
Layer1+2 与\textbf{完整版}恢复 Q1 均值（F1 上约 $6.3\times10^6$）；\textbf{Layer2（相对阈值）必要，Layer3（相位采样）边际}。

%==============================================================================
\section{实验 17：相位采样与 $\eta$ 策略消融（F1--F6）}
%==============================================================================
\subsection{实验说明}
在 Layer2 已启用前提下，比较\textbf{关闭相位采样（NoPhase）} vs \textbf{完整版相位采样}，以及四种 $\eta$（惩罚更新步长）策略。

\subsection{实验计算公式}
$\eta$ 为惩罚更新步长/比率；策略 Fixed / Phase / Eta\_RL / Unified 仅改变 $\alpha_i$ 更新规则，门控不变。

\subsection{算法伪代码}
固定 Layer1+2+变量选择，切换 Phase 或 $\eta$ 模块，25-run $\times$ F1--F6。

\subsection{结果评估指标公式}
final $F$ 均值$\pm$标准差；期望：各配置 endpoint 接近（差异 $<5\%$）。

\subsection{实验设置与参数}
25 runs，LLSO；详见 Table \texttt{tab:phase\_eta}（en\_ready 版）。

\subsection{实验结果}
NoPhase 与\textbf{完整版}、四种 $\eta$ 策略在 F1--F6 上 endpoint 几乎相同；主要影响通信时机，非最终质量。

%==============================================================================
\section{实验 18：深度 RL 架构对比（F1--F6）}
%==============================================================================
\subsection{实验说明}
替换默认 PPO 头为 SimpleNet/A2C/LSTM/MLP/DQN/Dueling DQN，检验 RL \textbf{架构}是否影响 endpoint。

\subsection{实验计算公式}
\begin{equation}
  \alpha_i^{t+1} = \pi_\theta(\mathrm{state}_i^t),\quad \pi_\theta \in \{\text{PPO, A2C, DQN, \ldots}\}
\end{equation}

\subsection{算法伪代码}
固定门控与 eva 上限，仅替换 $\pi_\theta$ 网络与训练算法，25-run $\times$ F1--F6。

\subsection{结果评估指标公式}
final $F$ 均值/标准差；架构间相对差异 $<5\%$ 视为等价。

\subsection{实验设置与参数}
25 runs，LLSO；Exp6 DeepRL 目录。

\subsection{实验结果}
\tryfig{media/fig_deeprl_comparison.pdf}{六种 RL 架构 F1--F6 最终 fitness 对比}
PPO 非唯一最优；各架构差异小于 5\%，与实验 15 结论一致。

%==============================================================================
\section{实验 19：RL 轨迹诊断（$\alpha$, $\rho$, conflict, reward）}
%==============================================================================
\subsection{实验说明}
25-run 均值轨迹展示 PPO 默认控制器\textbf{调整了什么}，非 RL 优于启发式的证据。

\subsection{实验计算公式}
\begin{equation}
  \mathrm{reward}_t = f(\Delta f_{\mathrm{pure}}, \mathrm{CI}^t, \alpha^t),\quad
  \rho^t = \alpha^t / \alpha^{\mathrm{base}}
\end{equation}

\subsection{算法伪代码}
聚合 25 条轨迹，按累计 eva 对齐后计算 $\alpha, \rho, \mathrm{CI}, \mathrm{reward}$ 的 run-mean $\pm$ std。

\subsection{结果评估指标公式}
时序曲线 + 跨 run 标准差带；无显著性检验。

\subsection{实验设置与参数}
F1--F6；日志 \texttt{RL-output\_runs25\_rho/}；脚本 \texttt{plot\_rl\_metrics\_runs25\_by\_metric.py}。

\subsection{实验结果}
\tryfig[0.48\textwidth]{media/rl_metrics_mean25_rho_alpha.pdf}{25-run 均值：$\alpha$ 轨迹}
\tryfig[0.48\textwidth]{media/rl_metrics_mean25_rho_rho.pdf}{25-run 均值：$\rho$ 轨迹}
\tryfig[0.48\textwidth]{media/rl_metrics_mean25_rho_conflict.pdf}{25-run 均值：冲突代理轨迹}
\tryfig[0.48\textwidth]{media/rl_metrics_mean25_rho_reward.pdf}{25-run 均值：RL reward 轨迹}

%==============================================================================
\section{实验 20：冲突--$\alpha$ 分箱关系}
%==============================================================================
\subsection{实验说明}
汇总 F1--F6 全部 run，按冲突 CI 分箱展示 $\alpha$ 的响应模式。

\subsection{实验计算公式}
\begin{equation}
  \bar{\alpha}_q = \text{均值}\bigl\{\alpha^t \mid \mathrm{CI}^t \in \mathrm{bin}_q\bigr\}
\end{equation}
即 CI 落在第 $q$ 箱时，惩罚系数 $\alpha$ 的样本均值（原写法 $\mathbb{E}[\alpha \mid \mathrm{CI} \in \mathrm{bin}_q]$ 同义）。

\subsection{算法伪代码}
脚本 \texttt{plot\_conflict\_alpha\_bins.py}；输入 25-run 轨迹。

\subsection{结果评估指标公式}
各箱 $\bar{\alpha}_q$（条件均值）；横轴为 CI 分箱，纵轴为 $\alpha$。

\subsection{实验设置与参数}
F1--F6 pooled；25 runs。

\subsection{实验结果}
\tryfig{media/conflict_alpha_bins_F1_F6.pdf}{F1--F6 冲突 CI 分箱 vs $\alpha$ 响应}

%==============================================================================
\section{实验 21：逐循环归一化漂移验证}
%==============================================================================
\subsection{实验说明}
为 Proposition fail-safe 有界漂移假设提供\textbf{经验支撑}，非 MACPO 收敛证明。

\subsection{实验计算公式}
\begin{align}
  e_d^t &= \frac{|x_d^t - r_d^t|}{R_d}
  \quad\text{（$e_d^t$：共享维 $d$ 的归一化不一致度）} \nonumber\\
  \varepsilon &= |e_d^{t+1} - e_d^t|
  \quad\text{（$\varepsilon$：静默外循环上单步漂移）} \nonumber\\
  e_d^{t+\tau} &\le e_d^t + (K-1)\varepsilon_{\max},\quad \tau \le K-1
  \quad\text{（fail-safe 命题上界）}
\end{align}

\subsection{算法伪代码}
从日志末尾读取 \texttt{diff\_mean}、\texttt{diff\_p95}、\texttt{diff\_max}（单步漂移 $\varepsilon$ 的均值/95 分位/最大值），跨 25 run 汇总。
\varnote{
\texttt{diff\_mean}：$\varepsilon$ 的 run 内均值；\texttt{diff\_p95}：95 分位；\texttt{diff\_max}：最大值。
}

\subsection{结果评估指标公式}
Mean $\bar\varepsilon$，Max $\varepsilon$，p95 $\varepsilon$（F1--F6 pooled）。

\subsection{实验设置与参数}
F1--F6 headline \textbf{完整版 RL-MACPO}，25 runs/函数。

\subsection{实验结果}
\begin{center}\small
\begin{tabular}{@{}lrrr@{}}
\toprule
函数 & Mean $\bar\varepsilon$ & Max $\varepsilon$ & p95 $\varepsilon$ \\
\midrule
F1 & 0.00409 & 0.0105 & 0.00655 \\
F2 & 0.00358 & 0.00727 & 0.00544 \\
F3 & 0.00412 & 0.00673 & 0.00566 \\
F4--F6 & 0.0037--0.0046 & $<0.01$ & $<0.008$ \\
\bottomrule
\end{tabular}
\end{center}
漂移处于窄带，支持 fail-safe 工程假设。

%==============================================================================
\section{实验 22：多区域经济调度 MAED-13}
%==============================================================================
\subsection{实验说明}
13 台机组、1800\,MW 总负荷的多区域经济调度；验证门控在\textbf{真实成本结构}下的通信削减与成本保持。

\subsection{实验计算公式}
\begin{equation}
  \min \; f_{\mathrm{pure}} = \sum_{i} C_i(P_i),\quad
  C_i(P_i) = a_i P_i^2 + b_i P_i + c_i
\end{equation}
subject to 区域功率平衡与联络线一致性（共享变量约束，MACPO 协商）。

\subsection{算法伪代码}
\begin{algorithm}[H]
\caption{MAED-13 paired dispatch runs}
\begin{algorithmic}[1]
\For{$\textit{run\_id} = 1$ \textbf{to} $25$}
  \State $\textsc{PairSeed} \gets \textit{run\_id}$
  \State Run baseline MACPO ($g^t \equiv 1$) until \textit{eva\_limit}
  \State Run full RL-MACPO with conflict gate until the same \textit{eva\_limit}
  \State Log best $f_{\mathrm{pure}}$ and $\bar p_{\mathrm{comm}}$
\EndFor
\end{algorithmic}
\end{algorithm}
\varnote{
$f_{\mathrm{pure}}=\sum_i C_i(P_i)$：纯发电成本（不含惩罚）。
$\bar p_{\mathrm{comm}}$：外循环协商触发率。
\textit{eva\_limit}：与 MACPO 相同的评估次数上限。
}

\subsection{结果评估指标公式}
\begin{equation}
  \Delta = \frac{f_{\mathrm{pure}}^{\mathrm{MACPO}} - f_{\mathrm{pure}}^{\mathrm{RL-MACPO}}}{|f_{\mathrm{pure}}^{\mathrm{MACPO}}|}\times 100\%
\end{equation}
$\Delta$：RL-MACPO 相对 MACPO 的纯目标改善百分比；$\Delta>0$ 表示 RL-MACPO 更优（$f_{\mathrm{pure}}$ 更低）。EV 场景符号约定见实验 24。

\subsection{实验设置与参数}
25 paired runs；\texttt{run\_maed.sh}；输出 \texttt{maed\_20260622\_114525/}。

\subsection{实验结果}
RL 通信 11.3\%；MACPO $f_{\mathrm{pure}} = 3.42\times10^4 \pm 38$；RL $3.43\times10^4 \pm 105$；$\Delta \approx -0.05\%$；通信削减 88.7\%。

%==============================================================================
\section{实验 23：资源-constrained 分布式调度}
%==============================================================================
\subsection{实验说明}
资源约束下的分布式任务/资源分配仿真；共享变量为跨节点资源配额。

\subsection{实验计算公式}
\begin{equation}
  \min \; f_{\mathrm{pure}} = \sum_i \bigl(\text{delay}_i + \text{cost}_i\bigr),\quad
  \text{s.t. 共享资源一致性}
\end{equation}

\subsection{算法伪代码}
同实验 22 配对流程；场景脚本 \texttt{run\_paper\_scenarios.sh}（RESOURCE）。

\subsection{结果评估指标公式}
同式 $\Delta$；25-run 均值$\pm$标准差。

\subsection{实验设置与参数}
25 paired runs；输出 \texttt{paper\_20260622\_114518/RESOURCE/}。

\subsection{实验结果}
RL 通信 11.6\%；MACPO $1.33\times10^3 \pm 184$；RL $\mathbf{1.18\times10^3 \pm 167}$；$\Delta = +11.0\%$（RL 更优）；通信削减 88.4\%。

%==============================================================================
\section{实验 24：电动汽车充放电协调（EV Dispatch）}
%==============================================================================
\subsection{实验说明}
EV 集群充放电与电网/站点协调；$f_{\mathrm{pure}}$ 为\textbf{净收益}（含售电收益），越负越好。

\subsection{实验计算公式}
\begin{equation}
  f_{\mathrm{pure}} = f_{\mathrm{cost}} - f_{\mathrm{gain}},\quad
  \text{更负的 } f_{\mathrm{pure}} \Rightarrow \text{更高净收益}
\end{equation}

\subsection{算法伪代码}
同实验 22；场景 EVDISPATCH。

\subsection{结果评估指标公式}
\begin{equation}
  \Delta = \frac{f_{\mathrm{pure}}^{\mathrm{MACPO}} - f_{\mathrm{pure}}^{\mathrm{RL}}}{|f_{\mathrm{pure}}^{\mathrm{MACPO}}|}\times 100\%
\end{equation}
正值表示 RL 净收益更高（$f_{\mathrm{pure}}$ 更负）。

\subsection{实验设置与参数}
25 paired runs；同 \texttt{paper\_20260622\_114518/EVDISPATCH/}。

\subsection{实验结果}
RL 通信 11.6\%；MACPO $-3.203 \pm 0.579$；RL $\mathbf{-3.534 \pm 0.987}$；$\Delta = +10.3\%$；通信削减 88.4\%。

%==============================================================================
\section{实验 25：IEEE 30/57/118 输电网区域调度}
%==============================================================================
\subsection{实验说明}
IEEE 标准测试系统上的\textbf{分区经济调度}；fail-safe $K{=}2$；检验大规模电网上的通信--成本权衡。

\subsection{实验计算公式}
\begin{equation}
  f_{\mathrm{pure}} = \sum_{\mathrm{regions}} \sum_{g \in \mathrm{region}} C_g(P_g),\quad
  \text{共享变量：联络线/区域边界功率}
\end{equation}

\subsection{算法伪代码}
\begin{algorithm}[H]
\caption{IEEE regional dispatch paired runs}
\begin{algorithmic}[1]
\For{$\textit{run\_id} = 1$ \textbf{to} $25$}
  \State Run MACPO and full RL-MACPO with fail-safe $K=2$
  \State Log best-so-far $f_{\mathrm{pure}}$ and $\bar p_{\mathrm{comm}}$
  \State Compute paired $\Delta_k$; aggregate median and Wilcoxon $p$
  \State Mark success if $f_{\mathrm{pure}}^{\mathrm{RL}} \le 2 \times f_{\mathrm{pure}}^{\mathrm{MACPO}}$
\EndFor
\end{algorithmic}
\end{algorithm}
\varnote{
$\Delta_k$：第 $k$ 次 run 上 MACPO 相对 RL-MACPO 的成本变化百分比（见下式）。
Success：灾难检测——RL 纯成本不超过 MACPO 的 2 倍即视为未崩溃。
$K=2$：IEEE 试点 fail-safe 周期（主基准默认 $K=10$）。
}

\subsection{结果评估指标公式}
\begin{align}
  \Delta_k &= \frac{f_{\mathrm{pure},k}^{\mathrm{MACPO}} - f_{\mathrm{pure},k}^{\mathrm{RL-MACPO}}}{|f_{\mathrm{pure},k}^{\mathrm{MACPO}}|}\times 100\% \\
  \text{Median paired } \Delta &= \mathrm{median}_k(\Delta_k) \quad\text{（25 次配对 run 的中位改善）} \\
  \text{Wilcoxon} &: \text{对 } \{\Delta_k\} \text{ 做双侧符号秩检验}
\end{align}

\subsection{实验设置与参数}
25 paired runs；IEEE30/57/118；批次 \texttt{114805/114810/113148}；\texttt{run\_power.sh}。

\subsection{实验结果}
\begin{center}\scriptsize
\begin{tabular}{@{}lrcrrc@{}}
\toprule
案例 & Succ. & RL comm. & Median $\Delta$ & Wilcoxon $p$ & Comm. drop \\
\midrule
IEEE-30 & 23/25 & 58.6\% & +0.5\% & 0.94 & $-$41.4\% \\
IEEE-57 & 25/25 & 58.8\% & $-$0.0\% & 0.67 & $-$41.2\% \\
IEEE-118 & 25/25 & 40.0\% & $-$0.6\% & 0.01 & $-$60.0\% \\
\bottomrule
\end{tabular}
\end{center}
IEEE-30/57 中位成本近零；118 节点有小幅可检测成本偏移（精度--通信权衡）；30 有两例 LLSO 停滞离群。

%==============================================================================
\section{实验 26：可扩展性 pilot（F1/F7/F13 + F1S50/S100）}
%==============================================================================
\subsection{实验说明}
随智能体数量增加（20/40/60/50/100），观察通信率与 fitness 变化；\textbf{pilot 规模小于主表}。

\subsection{实验计算公式}
同 Section~\ref{sec:notation} 式\eqref{eq:global}--\eqref{eq:gate}；F1S50/S100 为 F1 的 50/100 智能体扩展实例。

\subsection{算法伪代码}
F1/F7/F13：5-run pilot；F1S50/S100：25-run；MACPO vs RL-MACPO 配对。

\subsection{结果评估指标公式}
$\bar p_{\mathrm{comm}}$ + final fitness 均值$\pm$标准差。

\subsection{实验设置与参数}
5-run（F1/F7/F13）或 25-run（F1S50/S100）；\textbf{不可与 25-run 主表直接混读}。

\subsection{实验结果}
\begin{center}\small
\begin{tabular}{@{}lccccc@{}}
\toprule
基准 & 智能体 & MACPO comm. & RL comm. & MACPO $F$ & RL $F$ \\
\midrule
F1 & 20 & 100\% & 18.3\% & $1.88\times10^7$ & $4.95\times10^7$ \\
F7 & 40 & 100\% & 8.0\% & $7.24\times10^7$ & $9.54\times10^7$ \\
F13 & 60 & 100\% & 8.0\% & $1.17\times10^8$ & $1.15\times10^8$ \\
F1S50 & 50 & 100\% & 24.3\% & $1.69\times10^7$ & $2.09\times10^7$ \\
F1S100 & 100 & 100\% & 23.3\% & $1.47\times10^7$ & $1.95\times10^7$ \\
\bottomrule
\end{tabular}
\end{center}
5-run pilot 方差较大；F13 RL 略优于 MACPO；F1S50/S100 需更大 runs 才能与主表同等置信。

%==============================================================================
\section*{附录：应用案例 run 级 $\Delta$ 分布（实验 22--25 补充）}
\addcontentsline{toc}{section}{附录：应用 run 级分布}
%==============================================================================
\begin{center}\small
\begin{tabular}{@{}lrrr@{}}
\toprule
场景 & Median $\Delta$(\%) & Q1--Q3(\%) & Min--Max(\%) \\
\midrule
MAED-13 & 0.0 & [0.0, 0.0] & [$-$1.5, +0.4] \\
Resource scheduling & +9.5 & [$-$0.1, +19.5] & [$-$25.8, +42.9] \\
EV dispatch & +0.6 & [$-$1.2, +3.4] & [$-$4.2, +88.7] \\
IEEE-30 & +0.5 & [$-$2.8, +2.1] & [$-$823.7, +7.1] \\
IEEE-57 & 0.0 & [$-$6.5, +4.8] & [$-$22.0, +20.2] \\
IEEE-118 & $-$0.6 & [$-$1.0, +0.1] & [$-$4.2, +1.2] \\
\bottomrule
\end{tabular}
\end{center}

\vspace{1em}
\noindent\textit{编译说明}：使用 \texttt{xelatex experiment\_summary\_zh.tex}（两遍以生成目录）。若需嵌入结果图，请先运行 \texttt{scripts/plot\_*.py} 生成 \texttt{media/*.pdf}。

\end{document}
